% ==================================================
% Summation by Parts
% Discrete Calculus
% ==================================================

\subsection*{Summation by Parts}

\begin{exposition}
Summation by parts is the discrete analogue of integration by parts.
It is derived from the product rule exactly as the continuous version
is derived from the Leibniz rule.
\end{exposition}

\begin{theorembox}{Theorem (Summation by Parts)}
\begin{theorem}[Summation by Parts]
\label{thm:summation-by-parts}
For functions \(f, g : \mathbb{Z} \to \mathbb{R}\) and integers \(a \le b\),
\[
\sum_{k=a}^{b-1} f(k)\,\Delta g(k)
=
\Bigl[f(k)g(k)\Bigr]_{k=a}^{k=b}
-
\sum_{k=a}^{b-1} g(k+1)\,\Delta f(k).
\]
That is,
\[
\sum_{k=a}^{b-1} f(k)\,\Delta g(k)
=
f(b)g(b)-f(a)g(a)
-
\sum_{k=a}^{b-1} g(k+1)\,\Delta f(k).
\]
\hyperref[prf:summation-by-parts]{\textit{Go to proof.}}
\end{theorem}
\end{theorembox}

\begin{remark*}[Standard quantified statement]
\(\forall f,g : \mathbb{Z}\to\mathbb{R},\; \forall a,b\in\mathbb{Z},\; a\le b:\;
\displaystyle\sum_{k=a}^{b-1}f(k)\Delta g(k)
= f(b)g(b) - f(a)g(a) - \sum_{k=a}^{b-1}g(k+1)\Delta f(k)\).
\end{remark*}

\begin{remark*}[Predicate reading]
\[
\sum_{k=a}^{b-1}f(k)\Delta g(k)
=f(b)g(b)-f(a)g(a)-\sum_{k=a}^{b-1}g(k+1)\Delta f(k).
\]
\end{remark*}

\begin{remark*}[Interpretation]
The continuous integration by parts formula,
\[
\int_a^b f\,dg = \bigl[fg\bigr]_a^b - \int_a^b g\,df,
\]
is recovered from summation by parts by replacing \(\sum\) with
\(\int\) and \(\Delta\) with \(d\). The asymmetry in the discrete
formula --- \(g(k+1)\) rather than \(g(k)\) in the remaining sum ---
reflects the shift inherent in the product rule.
Summation by parts is useful for evaluating sums involving products,
particularly when one factor is easy to sum (as a falling factorial)
and the other is easy to difference. Classic applications include
evaluating \(\sum k r^k\) (where \(f(k) = k\) and
\(\Delta g(k) = r^k\), leading to a geometric-type sum) and establishing
Abel's summation formula.
\end{remark*}

\begin{dependencies}
\begin{itemize}
  \item \hyperref[thm:discrete-product-rule]{Discrete Product Rule}
  \item \hyperref[thm:fundamental-theorem-finite-calculus]{Fundamental Theorem of Finite Calculus}
\end{itemize}
\end{dependencies}
